\documentclass[10pt,a4paper]{amsart}
\usepackage[margin=19mm]{geometry}
\usepackage{amsmath,amssymb,mathtools}
\usepackage[scr=rsfs,cal=euler]{mathalfa}
\usepackage{xfrac}
\usepackage[unicode,hidelinks,bookmarks=false]{hyperref}
\newcommand{\ie}{i.e.\ }
\newcommand{\cf}{cf.\ }
\newcommand{\resp}{respectively\ }
\newcommand{\Subsection}{\subsection}
\providecommand{\nobreakdash}{\nobreak\hbox{-}}
\setlength{\emergencystretch}{2em}
\multlinegap0pt
\usepackage{amssymb}
\usepackage{mathtools}
\usepackage{booktabs}
\usepackage{caption}
\usepackage{tikz}
\usepackage{algorithm}\usepackage{algpseudocode}
\usepackage{xfrac}
\usepackage{float}
\usepackage[shortlabels]{enumitem}
\usepackage{bm}
\usepackage{xcolor}
\newtheorem{lemma}{Lemma}
\newtheorem{theorem}{Theorem}
\newtheorem{proposition}{Proposition}
\newcommand{\E}{\mathbb{E}}
\newcommand{\R}{\mathbb{R}}
\renewcommand{\d}{\mathrm{d}}
\renewcommand{\leq}{\leqslant}
\newcommand{\opnorm}[1]{{\left\| #1 \right\|}_{\mathcal{L}}}
\title{Submanifolds of class $C^{1,\alpha}$ and sets with positive $\mu$-reach}
\author{Vincent Borrelli, Jean-Baptiste Follet, Boris Thibert}
\date{Source: arXiv:2602.12999v1 (CC BY 4.0)}
\begin{document}
\maketitle

\noindent\emph{Source and licence.} Submanifolds of class $C^{1,\alpha}$ and sets with positive $\mu$-reach. Version arXiv:2602.12999v1 (CC BY 4.0), distributed by arXiv under the Creative Commons Attribution 4.0 International licence, http://creativecommons.org/licenses/by/4.0/, as recorded in the arXiv licence metadata for this exact version and shown on the abstract page https://arxiv.org/abs/2602.12999v1. Full licence notice: \texttt{licenses/arxiv-2602.12999-CC-BY-4.0.txt}.\par\smallskip

\noindent\textit{Editorial scope.} Counted unit: the article's proposition graph\_constant\_holder\_global (source0.tex lines 835--843). The article's complete original proof of it is delivered with it (source0.tex lines 845--1017, one \texttt{proof} environment of 173 lines). No mathematics is written, repaired or re-derived: the statement and the proof are the article's own lines, in the article's own notation, and no retained line is substituted or re-flowed.  Retained uncounted as the source-local context the proof needs: lines 583--595; lines 708--710; lines 754--763; lines 757--759.  Nothing was omitted from any retained span, so the package carries no omission record.  No retained line contains a citation, so the wrapper carries no bibliography and no bibliography entry.  No retained line contains a figure environment, an image-inclusion command or drawing code, so no image is delivered and the package declares no asset dependency.  Editorial formatting only, in addition: an AMS \texttt{amsart} preamble that restates the article's own declarations and macros used by the retained lines, namely the environments \texttt{lemma}, \texttt{theorem}, \texttt{proposition} on separate counters, together with the standard packages those lines need.  The retained environments print in this document as Lemma 1, Theorem 1, Proposition 1 in order of appearance, whereas the article prints the same items with the numbering of its own full document.  The complete native source of the article as distributed in the arXiv e-print for this exact version is bundled unchanged as \texttt{sources/194602/source0.tex}.
\begin{lemma} \label{lem: carac_c_1_alpha}
    Let $g : U \to F$ be a $C^{1}$ function, $\alpha \in [0,1]$ and $U$ a convex open set.
    The following statements are equivalent
    \begin{enumerate}
        \item there exists $C > 0$ such that for all $x,y \in U$,
    \[\opnorm{ \d f_x - \d f_y} \leq C \|x -y\|_E^\alpha,\]
    where $\opnorm{\cdot}$ denotes the operator norm.

        \item there exists $C' > 0$ such that for all $x,y \in U$,
    \[\| f(y) - f(x) - \d f_x(y-x)\|_F \leq C' \|y-x\|_E^{1 + \alpha}.\]
    \end{enumerate}
    Moreover, we have $(1+\alpha)\,  C \leq C' \leq 6 C$. Furthermore, the implication $(2) \implies (1)$ does not require the convex hypothesis on $U$.
\end{lemma}

\begin{theorem}[Inverse Function Theorem] \label{theo: inverse_function}
    Let $f : U \to F$ of class $C^{1,\alpha}$ with $\alpha \in [ 0 , 1]$ and let $a \in U$. We suppose that $\d f_a$ is an isomorphism of $\mathscr{L}(E,F)$. Then there exists an open set $V \subset U$ such that $a \in V$ and $f \vert_V$ is a $C^{1,\alpha}$-diffeomorphism on its image.
\end{theorem}

\begin{lemma} \label{lem: constant_holder_inversion_local}
    Let $V$ be an open set of $E$ and $f : V \to F$ be a $C^{1}$-diffeomorphism on its image.
    Let $W \subset V$ be a convex neighborhood of some $a \in V$ such that $\d f |_W$ is $\alpha$-Holder with $\alpha \in [ 0 , 1]$ and such that for all $x \in W$, 
    \begin{equation} \label{eq: condition_alpha_holder_sur_tout_ouvert}
        \opnorm{(\d f_x)^{-1}} \leq 2  \opnorm{(\d f_a)^{-1}} ~~ \text{and} ~~ \opnorm{ (\d f_a)^{-1} \circ \d f_x   - id_E } \leq \frac{1}{2}  .
    \end{equation}
    Then $\d f^{-1}$ is $\alpha$-Hölder on $f(W)$ and its Hölder constant $C'$ satisfies
    $$ C' \leq \frac{ 3.2^{3 + \alpha} \, \opnorm{ (\d f_a)^{-1} }^{2+\alpha}}{1 + \alpha} C$$
    where $C$ is the Hölder constant of $\d f|_W$.
\end{lemma}

    \begin{equation} 
        \opnorm{(\d f_x)^{-1}} \leq 2  \opnorm{(\d f_a)^{-1}} ~~ \text{and} ~~ \opnorm{ (\d f_a)^{-1} \circ \d f_x   - id_E } \leq \frac{1}{2}  .
    \end{equation}

\begin{proposition} \label{prop: graph_constant_holder_global}
    Let $M$ be a $C^{1, \alpha}$ submanifold of $\E^n$ and $q_0 \in M$.
    There exist $U \subset \E^n$ an open neighborhood of $q_0$ and $K  = K(U) > 0$ such that for all $q \in  M \cap U$
    \begin{itemize}
        \item[(1)] there is an open set $V_q \subset T_q M$ such that $M$ is locally the graph of a $C^{1, \alpha}$ function $h_{q} : V_q \subset T_{q}M \to N_{q}M$, \textit{i.e.} 
        $$ M \cap U = \big\{ \big(x , h_{q}(x)\big) ~| ~ x \in V_q  \subset T_q M\big\}, $$
        \item[(2)] $\d (h_{q})$ is $\alpha$-Hölder on $U$ and its Hölder constant is bounded by $K$.
    \end{itemize}
\end{proposition}

\begin{proof}
    For any $q \in M$, we denote by $\pi_{q}^T$ the orthogonal projection on $T_q M + q$ and  $\pi_{q}^N$ the orthogonal projection on $N_q M + q$. In order to prove the proposition, we just need to prove that
    there exist $U \subset \E^n$ an open neighborhood of $q_0$ and $K > 0$ such that 
    \begin{itemize}
        \item[(i)] for all $q \in  M \cap U$, $\pi_{q}^T |_{M\cap U}$ is a $C^{1,\alpha}$-diffeomorphism on its image,
        \item[(ii)] for all $q \in  M \cap U$, $\d \big( \pi_{q}^T |_{M\cap U} \big)^{-1}$ is $\alpha$-Hölder on $M \cap U$ and its Hölder constant is bounded by $K$.
    \end{itemize}
    Indeed, if we set $h_{q} = \pi_{q}^N \circ \big( \pi_{q}^T \big)^{-1}$ we obtain exactly $(1)$ and $V_q  = \pi_{q}^T(M\cap U)$.
    Then, since $\| \d (\pi^N_{q})\| \leq 1$ and by the chain rule, we have that $\d (h_{q})$ is $\alpha$-Hölder on $V_{q}$ with its Hölder constant bounded by $K$, hence $(2)$ is proved.
    
    The proof of $(i)$ and $(ii)$ is done in the next five steps.
    The first four steps will prove $(i)$ by mimicking the proof of the Theorem \ref{theo: inverse_function} and adding a parameter $q\in M.$
    We fix $q_0 \in M$ and $\psi : V_0 \subset \R^m \to U_0 \cap M$ a $C^{1,\alpha}$ parametrization of $M$ such that $\psi(0) = q_0$ where $U_0$ (resp. $V_0$) is an open sets of $\E^n$ (resp. $\R^m$).
    We also suppose $U_0$ small enough so that $\d \psi$ and $\d (\psi^{-1})$ are $\alpha$-Hölder.
    We denote by $C_{\psi}$ and $C_{\psi^{-1}}$ their Hölder constant on $U_0$ (resp. $V_0$).
    We define, for all $q \in U_0$
    \[ \begin{array}{ccccc}
    f_q & : & V_0 \subset \R^m & \longrightarrow & \R^m \\
     & & x & \longmapsto & \d (\psi^{-1})_{q} \circ t_{-q}  \circ \pi_{q}^T \circ \psi(x) \\
    \end{array}\]
    with $t_{-q} : x \mapsto x - q$.
    We also define $\varphi_{q} = f_{q} - id$ and
    for all $y \in \R^m$, we define $T_{y,q}(x) = y -\varphi_{q}(x)$. 
    We then get,
    $$ y = f_{q}(x) \iff T_{y,q}(x) = x.$$
    In order to show that $\pi_q^T |_M$ is locally a bijection, we need to show that $f_q$ is a bijection, which is recast as a fixed point problem for $T_{y,q}$\\
    
    \noindent \textbf{Step 1.} The goal of this step is to show that there exists $r > 0$ such that $T_{y,q}$ is $\frac{1}{2}$-Lipchitz on $B(0,r)$, for every $y \in \R^m$ and $q \in \psi \big( B(0,r) \big)$.    
    Using chain rules, we have for $x \in V_0$,
    \[ \d (\varphi_{q})_{x} = \d ( f_{q})_{x} - id = \d (\psi^{-1})_{q} \circ \d \big( \pi_{q}^T\big) \circ \d \psi_{x} - id .  \]
    Also, since $ \pi_{q}^T(q) = q$ and $\d (\pi_{q}^T)_{q} |_{T_q M} = id_{T_{q}M}$, we have
    $$\d (f_{q})_{\psi^{-1}(q)} = \d (\psi^{-1})_{q}  \circ \d (t_{-q}) \circ \d \big( \pi_{q}^T\big) \circ \d \psi_{\psi^{-1}(q)} = \d (\psi^{-1})_{q} \circ \d \psi_{\psi^{-1}(q)} = id_{\R^m}$$
    thus $\| \d (\varphi_{q_0})_{0} \| = 0$.
    Moreover, we know that  and that the application 
    \[ (x,q) \in V_0 \times M \cap U_0 \mapsto  \| \d (\psi^{-1})_{q} \circ \d \big( \pi_{q}^T\big) \circ \d \psi_{x} - id  \| \]
    is continuous.
    Thus, there exists $r > 0$ such that $B(0,r) \subset V_0$ and for all $(x, q) \in B(0,r) \times \psi \big( B(0,r) \big)$, we have
    \begin{equation} \label{eq: maitrise_condition_projection}
        \opnorm{ \d ( \varphi_{q})_{x} }  \leq \frac{1}{2}.
    \end{equation}
    Using the Mean value inequality, we get that for all $q \in \psi \big( B(0,r) \big)$ and $ x_1,x_2 \in  B(0,r)$,
    $$ \| \varphi_{q}(x_1) - \varphi_{q}(x_2) \| < \frac{1}{2} \| x_1 - x_2\|$$
    thus $$ \|T_{y,q}(x_1) -T_{y,q}(x_2) \| < \frac{1}{2} \| x_1 - x_2\|.$$

    \noindent \textbf{Step 2.}
    There exists $r' > 0$ such that 
    $$M \cap B(q_0,r') \subset \psi \big( B(0,\frac{r}{8}) \big)$$ 
    and we define
    $$U := B\Big(q_0, \frac{r'}{2}\Big).$$
    We fix $q \in U \cap M$. 
    Our goal is to prove that, up to a restriction to a well chose set, $\pi^T_{q}$ is a bijection on its image.
    We have that 
    $$U \cap M \subset M \cap B(q_0,r') \subset \psi \big( B(0,\frac{r}{8}) \big) \subset \psi \big( B(0,r) \big)$$
    thus for $q \in U \cap M$, $T_{y,q}$ is $\frac{1}{2}$-Lipchitz on $ B(0,r)$.
    Thus for $T_{y,q}$ to be a contraction, we just need to prove that $T_{y,q} \big( \overline{B(0,\frac{r}{2})} \big) \subset  B(0,\frac{r}{2})$.
    We now suppose that $y \in B(0,\frac{r}{8})$, then for all $x \in \overline{B(0,\frac{r}{2})}$,we have
    \begin{align*}
        \| \varphi_{q}(x) \| &= \| \varphi_{q}(x) - \varphi_{q}\big(\psi^{-1} (q)\big) \| \leq \frac{1}{2} \| x-\psi^{-1}(q) \| \\
        &\leq \frac{1}{2} \| x \| + \frac{1}{2} \| \psi^{-1}(q) \| < \frac{r}{4} + \frac{r}{8} \leq \frac{3r}{8}
    \end{align*}
    because $q \in U \cap M \subset \psi \big( B(0,\frac{r}{8}) \big)$.
    Thus, we get
    $$\|T_{y,q}(x) \| < \frac{3r}{8} + \|y \| <  \frac{3r}{8} + \frac{r}{8} \leq \frac{r}{2}.$$
    Since $  \overline{B(0,\frac{r}{2})}$ is a complete metric space, we can apply Banach fixed point theorem of $T_{y,q} |_{\overline{B(0,\frac{r}{2})}}$ for all $y \in B(0,\frac{r}{8})$.
    Hence, for all $y \in B(0,\frac{r}{8})$ there exists a unique $x \in B(0,\frac{r}{2})$ such that $f_q(x)=y$.
    Therefore, if we define $U_{q'} \subset \R^m$ such that 
    $$U_{q} = \big(f_{q}\big)^{-1} \big( B(0,\frac{r}{8}) \big)$$
    then $f_{q} |_{U_{q}}$ is a bijection on its image (being $B(0,\frac{r}{8})$).
    Since $\d (\psi^{-1})_q$ and $\psi$ are bijections, we have that 
    $\pi^T_{q} |_{\psi(U_{q})}$ is a bijection on its image for all $q \in U \cap M$.\\

    \noindent \textbf{Step 3.} The goal of this step is to prove $(i)$.
    We first show that, up to a smaller choice of $r'>0$, we have that for all $q \in M \cap B(q_0,r')$
    \begin{equation} \label{eq: inclusion_differentielle_parametrisation}
        T_{q} M \cap B(0,r') \subset  \d \psi_{\psi^{-1}(q)} \Big( B\big(0,\frac{r}{8}\big) \Big).
    \end{equation}
    Indeed, thanks to the continuity of $x \mapsto  \d \psi_{\psi^{-1}(x)}$, the function
    \[ \lambda : q \in  \psi\big(B\big(0,\frac{r}{8}\big)\big) \mapsto \underset{\|v\| = 1}{\sup} \|  \d \psi_{\psi^{-1}(q)} .v  \|\]
    is continuous. 
    Since it is defined on a compact set, it has a minimum $\lambda_0$.
    Then, we have for all $q \in \psi\big(B (0,\frac{r}{8})\big)$
    \[   T_{q} M \cap B(0,\frac{\lambda_0 r }{8}) \subset  T_{q} M \cap B(0,\frac{\lambda(q) \, r }{8}) \subset \d \psi_{\psi^{-1}(q)} \Big( B\big(0,\frac{r}{8}\big) \Big).\]
    Therefore, by taking $r'$ to be the minimum between the previous $r'$ and $\frac{\lambda_0 r }{8}$ we get \eqref{eq: inclusion_differentielle_parametrisation}.
    Now, we prove that
    \[ U \cap M \subset \underset{q \in  M \cap B(q_0,\frac{r'}{4})}{\bigcap} \psi(U_{q}).\]
    Let $q \in B\big(q_0, \frac{r'}{2}\big)$, since $\pi_q^T$ retracts distances, we have that
    $$  \pi^T_{q} \Big(  B\Big(q_0, \frac{r'}{2}\Big) \Big) \subset  B\Big( \pi_q^T(q_0), \frac{r'}{2}\Big)$$
    and
    $$ \| \pi_q^T (q_0) - q \| \leq \| q_0 - q \| \leq \frac{r'}{2} $$
    thus
    $$   t_{-q} \circ \pi^T_{q} \Big(  B\Big(q_0, \frac{r'}{2}\Big) \Big) \subset  B( 0 , r').$$
    Therefore,
    $$   t_{-q} \circ \pi^T_{q} ( U \cap M )  \subset B(0 , r') \cap T_q M.$$
    Thanks to Equation \eqref{eq: inclusion_differentielle_parametrisation}, we have
    \[   t_{-q} \circ \pi^T_{q} (U \cap M)  \subset \d \psi_{\psi^{-1}(q)} \big( B(0,\frac{r}{8}) \big) .\]
    Hence by the definition of $f_q$, we have than $ \big( t_q \circ \pi_q^T) ^{-1} \circ \d \psi_{\psi^{-1}(q)} = \psi \, \circ (f_q)^{-1}$, thus
    \[U \cap M \subset \psi \circ \big( f_{q} \big)^{-1} \Big( B\big(0,\frac{r}{8}\big) \Big)  = \psi (U_{q}).\]
    Now, thanks to Step 2, we have for all $q \in U \cap M$, $\pi_{q}^T |_U$ is a bijection on its image $V_{q}$. \\
    %{\color{red} \textbf{ The arguments of the proof of Theorem \ref{theo: inverse_function}, apply and show that the all the $\pi_{q}^T |_{U \cap M}$ are $C^{1,\alpha}$-diffeomorphisms on their respective image.}}\\

    \noindent \textbf{Step 4.}
    In this step, similarly as in the proof of Theorem \ref{theo: inverse_function}, we show that all the $\pi_{q}^T |_{U \cap M}$ are $C^{1,\alpha}$-diffeomorphisms on their respective image.
    
    \noindent We fix $q \in U \cap M$.
    Firstly, we will show that $f_q^{-1}$ is continuous. Let $y_1,y_2 \in B(0,\frac{r}{8})$. Then there exist $x_1,x_2 \in U_q$ such that $f_q(x_1) = y_1$ and $f_q(x_2) = y_2$.
    We have
    \[ f_q^{-1}(y_2) - f_q^{-1}(y_1) = x_2 - x_1 = T_{y_2,q}(x_2) - T_{y_1,q}(x_1) = y_2 - \varphi_q(x_2) - y_1 + \varphi_q(x_1). \]
    Since $\varphi_q$ is $\frac{1}{2}$-Lipschitz, we get
    \[ \| f_q^{-1}(y_2) - f_q^{-1}(y_1) \| \leq \| y _2 - y_1 \| + \frac{1}{2} \| x_2 - x_1\| = \| y _2 - y_1 \| + \frac{1}{2} \| f_q^{-1}(y_2) - f_q^{-1}(y_1) \| \]
    thus
    \begin{equation*} 
        \| f_q^{-1}(y_2) - f_q^{-1}(y_1) \| \leq 2 \| y _2 - y_1 \|.
    \end{equation*}
    This shows that the map $f_q^{-1}$ is $2$-Lipschitz and thus continuous. \\

    \noindent Secondly, we are going to show that $f_q^{-1}$ is of class $C^{1,\alpha}$ using Lemma \ref{lem: carac_c_1_alpha}.
    Since $\d (f_{q_0})_0$ is invertible and $(q,x) \mapsto \d (f_q)_x$ is continuous, we can suppose $r$ small enough so that for all $q \in U$ and $x \in U_q$, $\d (f_q)_x$ is invertible.
    We recall that  $\d \psi$ and $\d (\psi^{-1})$ are $\alpha$-Hölder on $U \cap M$ and $\psi(U \cap M)$.
    For all $q \in U \cap M$, $d \pi_q^T$ is 1-Lipchitz thus $\d (f_{q})$ is $\alpha$-Hölder on $W :=\psi^{-1}(U \cap M)$.
    By shrinking $U$ if needed, we can assume that $W$ is convex. 
    We fix $q \in U\cap M$, let $x_1,x_2\in W$, $y_1=f_q(x_1),$ $y_2=f_q(x_2),$ 
    $$ y_2 - y_1 = \d (f_q)_{x_1}(x_2-x_1) + R_{x_1}(x_2)$$
    where the remainder satisfies $R_{x_1}(x_2) = O(\|x_2-x_1\|^{1+\alpha})$, see Lemma~\ref{lem: carac_c_1_alpha}. 
    Composing by $\big(\d (f_q)_{x_1}\big)^{-1}$, we obtain
    \begin{equation} \label{eq: _inv_loc_definition_C_1_alpha_2}
        (f_q)^{-1}(y_2) -  (f_q)^{-1}(y_1) = x_2-x_1 = \big(\d (f_q)_{x_1}\big)^{-1}(y_2 - y_1) - \big(\d (f_q)_{x_1}\big)^{-1}\big(R_{x_1}(x_2)\big).
    \end{equation}
    Since $(f_q)^{-1}$ is $2$-Lipschitz, we have that $\|x_2-x_1\| \leq 2 \|y_2-y_1\|$ thus
    \begin{align*}
        \frac{\| \big(\d (f_q)_{x_1}\big)^{-1}\big(R_{x_1}(x_2)\big) \|}{\|y_2-y_1\|^{1+\alpha}} &\leq 2^{1+\alpha} \frac{\| \big(\d (f_q)_{x_1}\big)^{-1}\big(R_{x_1}(x_2)\big) \|}{\|x_2-x_1\|^{1+\alpha}} \\ 
        &\leq 2^{1+\alpha} \bigg\| \big(\d (f_q)_{x_1}\big)^{-1}\bigg( \frac{R_{x_1}(x_2)}{\|x_2-x_1\|^{1+\alpha}} \bigg) \bigg\| = O_{x_2\to x_1}(1).
    \end{align*}
    We deduce that 
    \begin{equation} \label{eq: inverse_c_1_alpha_2}
        \big(\d (f_q)_{x_1}\big)^{-1}\big(R_{x_1}(x_2)\big) = O(\|y_2-y_1\|^{1+\alpha}).
    \end{equation}
    Thus by Equation \eqref{eq: _inv_loc_definition_C_1_alpha_2}, $f_q^{-1}$ is differentiable and 
    $$\d \big( (f_q)^{-1} \big)_{y_1} = \big(\d (f_q)_{(f_q)^{-1}(y_1)} \big)^{-1}.$$
    Since $f_q^{-1}$ and $\d (f_q)$ are continuous, $f_q^{-1}$ is of class $C^1$. Equation \eqref{eq: inverse_c_1_alpha_2} and Lemma \ref{lem: carac_c_1_alpha} imply that  $f_q^{-1}$ is of class $C^{1,\alpha}$.
    For all $q \in U$, we recall that
    $$(\pi_{q}^T)^{-1} |_{V_q} = \psi \circ f_{q}^{-1} \circ \d \big(\psi^{-1}\big) _q \circ t_{-q} |_{V_q} $$
    where $V_q = \pi_{q}^T (U \cap M)$.
    Thus the $\pi_{q}^T |_{U \cap M}$ are $C^{1,\alpha}$-diffeomorphisms on their respective image.\\
    
    
    \noindent \textbf{Step 5.} 
    In this last step, we will prove the point $(ii)$ of the proposition.
    %The proof of $(ii)$ reduces to apply Lemma \ref{lem: constant_holder_inversion_local} to the projections $\pi_{q}^T |_U$.
    %However since $U = B\big(q_0, \frac{r'}{2}\big) \cap M $ is not a Banach space, we consider instead the map $f_q$.
    We denote by $C_{\psi}$ and $C_{\psi^{-1}}$ the $\alpha$-Hölder of constant of $\d \psi$ and $\d (\psi^{-1})$ on $U \cap M$ and $\psi(U \cap M)$.
    For all $q \in U \cap M$, since $d \pi_q^T$ is 1-Lipchitz, we have that $\d(f_{q})$ is $\alpha$-Hölder on $W = \psi^{-1}(U \cap M)$, we denote its constant by $C_{f_{q}}$.
    We have 
    \[ C_{f_{q}} \leq \opnorm{\d \psi_{q}} \, C_{\psi} \leq \opnorm{(\d \psi)_{|U}  }\, C_{\psi}\]
    where the right term is finite since $U$ is bounded.
    Moreover, for all $q\in U \cap M$, since $\d (f_{q})_{\psi^{-1}(q)} = id_{\R^m}$ and using Equation \eqref{eq: maitrise_condition_projection}, we get for all $x \in W$
    \[ \opnorm{\d (f_q)_x - id_{\R^m} } = \opnorm{ \d ( \varphi_q)_x } \leq \frac{1}{2}\]
    which is exactly the first inequality of Condition \eqref{eq: condition_alpha_holder_sur_tout_ouvert} of Lemma \ref{lem: constant_holder_inversion_local} at the point $\psi^{-1}(q)$ and for the open set $W$.
    Moreover, we have
    \[ \opnorm{ \big( \d (f_q)_x \big)^{-1} } -1 \leq \opnorm{  \big( \d (f_q)_x \big)^{-1} - id_{\R^m} } \leq \opnorm{ \big( \d (f_q)_x \big)^{-1} }  \opnorm{ \d (f_q)_x - id_{\R^m}  } \]
    thus using the two previous inequality, we have
    \[ \opnorm{ \big( \d (f_q)_x \big)^{-1} } -1 \leq \frac{1}{2} \opnorm{ \big( \d (f_q)_x \big)^{-1} } \]
    which implies that for all $x \in W$,
    \[ \opnorm{ \big( \d (f_q)_x \big)^{-1} } \leq 2\]
    which is the second inequality of Condition \eqref{eq: condition_alpha_holder_sur_tout_ouvert}.
    Thus we can apply Lemma \ref{lem: constant_holder_inversion_local} to $f_{q}$ at the point $\psi^{-1}(q)$ and for the convex open set $W$.
    Thus $\d (f_{q})^{-1}$ is $\alpha$-Hölder on $f_{q}(W') = \d \psi_{q}(V_{q} - q)$ and its constant $C_{f^{-1}_{q}}$ verifies
    \[ C_{f^{-1}_{q}} \leq \frac{3.2^{3+\alpha} \opnorm{ \big(\d (f_{q})_{\psi^{-1}(q)} \big)^{-1} }}{1 + \alpha} C_{f_{q}} \leq \frac{3.2^{3+\alpha}}{1 + \alpha} \, {\opnorm{(\d \psi)_{|U} }} \, C_{\psi}\]
    since $\d (f_{q})_{\psi^{-1}(q)} = id_{\R^m}$.
    We recall that
    $$(\pi_{q}^T)^{-1} |_{V_q} = \psi \circ f_{q}^{-1} \circ d \big(\psi^{-1}\big) _q \circ t_{-q} |_{V_q} $$
    and that $\d (\psi^{-1})$, $d( f_{q}^{-1})$ are $\alpha$-Hölder on receptively $\psi(U \cap M)$ and $f_{q}(W)$.
    Thus $\d \big( (\pi_{q}^T)^{-1} \big) |_{U \cap M}$ is $\alpha$-Hölder on $V_{q}$ and its Hölder constant is bounded by a constant that does not depend on $q$.
\end{proof}

\end{document}
